% sections/method.tex -- Section 5 (paper_v2_vn, translated in L5 from paper_v2's
% sections/method.tex; terms as in glossary_vn.tex). Every number is \val{} from numbers.tex.
% Figure: figures/f2_ties.tex (F2). F1 (the overview), the changes from the base papers with their
% reasons and the search's details are in the supplementary document (supplement_vn.tex, N31). The
% key settings are in Section 6.
\section{Phương pháp}
\label{sec:method}

\subsection{Tìm kiếm meta-path cho dự~đoán~liên~kết}
\label{sec:search}

\paragraph{LMSPS\@.} Trước hết, LMSPS \citep{li2024longrange} cho mỗi meta-path ứng
viên một ma trận đặc trưng trên các nút đích: các đặc trưng gốc ở đầu kia của meta-path được lấy
trung bình qua từng bước dọc theo meta-path. Một siêu mạng gán cho mỗi ứng viên $k$ một trọng số
$\alpha_k$. Ở mỗi epoch siêu mạng này hợp nhất $M$ meta-path lấy mẫu từ một nhóm ứng viên với trọng số
Gumbel-softmax tính từ $\alpha$. Nhóm ứng viên thu hẹp dần còn $V$ meta-path có $\alpha$ lớn
nhất. Trọng số của mạng được học trên dữ liệu huấn luyện và $\alpha$ trên tập dữ liệu xác thực. Cuối cùng, các tập $M$ meta-path từ
nhóm ứng viên cuối được tính điểm trên tập dữ liệu xác thực. Tập tốt nhất là kết quả. Chúng tôi bắt
đầu từ mã nguồn của các tác giả\footnote{\url{https://github.com/JHL-HUST/LMSPS}}, vốn được xây
dựng cho bài toán phân loại nút. Chúng tôi thay đổi mã nguồn cho bài toán dự đoán liên kết.

\paragraph{Điểm và hàm mất mát.} Embedding $h_p$ của một sản phẩm là tổng có trọng số của các đặc
trưng meta-path của nó, mỗi đặc trưng qua một ánh xạ tuyến tính riêng, tiếp theo là PReLU và
dropout. Mỗi người dùng có một embedding tự do $e_u$ \val{search-hidden} chiều. Một cặp được tính điểm
bằng
\begin{equation}
  s(u, p) = e_u^{\top} h_p ,
  \label{eq:score}
\end{equation}
và được huấn luyện với hàm mất mát BPR \citep{rendle2009bpr}, $-\log \sigma\bigl(s(u, p^{+}) -
s(u, p^{-})\bigr)$, trên các cạnh huấn luyện $(u, p^{+})$ và các sản phẩm $p^{-}$ lấy mẫu ngẫu nhiên
từ những sản phẩm mà $u$ không có cạnh huấn luyện nối tới. HAN sẽ dùng cùng thang điểm và hàm mất mát này.

\paragraph{Ứng viên.} Các ứng viên là những meta-path dài từ 1 đến \val{search-hops} bước bắt đầu và
kết thúc tại một sản phẩm, vì HAN cần một đồ thị sản phẩm cho mỗi meta-path:
\val{search-cands-raw} meta-path, hay \val{search-cands} khi mỗi meta-path được gộp với meta-path
ngược của nó cho cùng một đồ thị sau khi đối xứng hóa.

\paragraph{Biện pháp chống rò rỉ.} Các đặc trưng được tính trên đồ thị huấn luyện, vốn chứa cả chính
các cạnh huấn luyện. Vì vậy, với một cạnh huấn luyện, đặc trưng meta-path qua người dùng của sản phẩm
đã chứa sẵn người dùng của cạnh đó, điều mà một cạnh giữ lại không bao giờ có. Để tránh điều này,
chúng tôi học $\alpha$ trên các cạnh xác thực. Ở bước cập nhật trọng số, chúng tôi tính đặc trưng
meta-path qua người dùng \val{search-folds} lần, mỗi lần bỏ đi một phần tư khác nhau của các cạnh
huấn luyện, rồi đọc đặc trưng của mỗi cạnh từ phiên bản không chứa cạnh đó.

\paragraph{Lịch huấn luyện và cách chọn.} Một tập gồm $M = 4$ hoặc $M = 6$ meta-path. Nhóm ứng
viên cuối gồm $V = \val{search-pool}$ meta-path, để có thể tính điểm cả \val{search-sets} tập của nhóm. Một lần chạy
kéo dài \val{search-epochs} epoch, \val{search-warmup} epoch đầu với toàn bộ nhóm ứng viên. LMSPS
tính điểm các tập lấy mẫu với nhiễu Gumbel mới, mà trên dữ liệu của chúng tôi làm điểm của một tập
thay đổi nhiều hơn chênh lệch giữa các tập. Chúng tôi tính điểm mỗi tập một lần, không có nhiễu,
bằng BPR xác thực của nó so với \val{split-negatives} mẫu âm cố định của mỗi cạnh xác thực. Mỗi lần
chạy cũng ghi lại \emph{nhóm đồng hạng} của nó, tức các tập nằm trong phạm vi
$\delta = \val{search-delta}$ so với tập có điểm tốt nhất.

\subsection{Một quyết định từ nhiều lần chạy}
\label{sec:fleet}

Một lần chạy không xếp hạng được các tập tốt nhất của nó (Hình~\ref{fig:ties}): nhóm đồng hạng
của nó chứa trung vị \val{search-tie-med-4} trong \val{search-sets} tập với $M = 4$ meta-path và
\val{search-tie-med-6} với $M = 6$ meta-path. Qua \val{search-runs} lần chạy cho mỗi kích thước, có
\val{search-winners-4} và \val{search-winners-6} tập thắng khác nhau. Vì vậy, chúng tôi chạy tìm
kiếm \val{search-runs} lần cho mỗi kích thước, với seed từ 1 đến \val{search-runs} trên cùng dữ liệu
và phép chia, rồi đưa ra một quyết định từ tất cả các lần chạy.

% f2_ties.tex -- F2: one run cannot rank its top sets (paper_v2_vn, translated in L3 from
% paper_v2's figures/f2_ties.tex). Data: f2_scores.dat and f2_picks.dat, written by
% make_numbers.py from the Fleet v3 records and re-selections.
\begin{figure}[htbp]
  \centering
  \begin{tikzpicture}
    \begin{axis}[
        name=scores, width=\linewidth, height=3.9cm,
        xlabel={Tập, tốt nhất trước}, ylabel={BPR xác thực\\vượt tập tốt nhất},
        ylabel style={align=center},
        xmin=0, xmax=211, ymin=0, ymax=0.125, xtick={1,50,100,150,210},
        yticklabel style={/pgf/number format/fixed, /pgf/number format/precision=2},
        scaled y ticks=false, tick label style={font=\scriptsize}, label style={font=\footnotesize},
        title={(a) Một lần chạy (seed 1): \val{search-sets} tập của nó},
        title style={font=\footnotesize},
        legend style={font=\scriptsize, at={(0.03,0.97)}, anchor=north west, draw=none,
          fill=white, fill opacity=0.85, text opacity=1},
        legend cell align=left, axis line style={black!60}, tick style={black!60}]
      \fill[cSearch!10] (axis cs:0,0) rectangle (axis cs:211,0.010);
      \addplot[cSearch, line width=1pt] table[x=rank, y=gapfour] {figures/f2_scores.dat};
      \addlegendentry{4 meta-path: \val{search-tie-seed1-4} tập trong phạm vi $\delta$}
      \addplot[cSearch!55, line width=1pt, densely dashed] table[x=rank, y=gapsix]
        {figures/f2_scores.dat};
      \addlegendentry{6 meta-path: \val{search-tie-seed1-6} tập trong phạm vi $\delta$}
      \draw[black!55, dashed] (axis cs:0,0.010) -- (axis cs:211,0.010)
        node[pos=0.03, above right, font=\scriptsize, text=black!75]
        {dung sai $\delta = \val{search-delta}$};
    \end{axis}
    \begin{axis}[
        at={(scores.below south west)}, anchor=above north west, yshift=-0.35cm,
        width=\linewidth, height=3.3cm,
        ybar, bar width=9pt, ymin=0, ymax=31, ylabel={Số tập khác nhau},
        symbolic x coords={4,6}, xtick=data, xticklabels={4 meta-path, 6 meta-path},
        enlarge x limits=0.5,
        tick label style={font=\scriptsize}, label style={font=\footnotesize},
        title={(b) Lựa chọn qua \val{search-runs} lần chạy}, title style={font=\footnotesize},
        legend style={font=\scriptsize, at={(0.5,-0.24)}, anchor=north, draw=none,
          legend columns=1}, legend cell align=left,
        axis line style={black!60}, tick style={black!60}]
      \addplot[fill=cSearch, draw=white] table[x=size, y=own] {figures/f2_picks.dat};
      \addlegendentry{tập thắng riêng của từng lần chạy}
      \addplot[fill=cSearch!45, draw=white] table[x=size, y=vote] {figures/f2_picks.dat};
      \addlegendentry{bỏ phiếu qua các lần chạy}
      \addplot[fill=cOther!45, draw=white, error bars/.cd, y dir=both, y explicit,
          error bar style={black!70, line width=0.6pt}]
        table[x=size, y=nullmed, y error plus=errhi, y error minus=errlo] {figures/f2_picks.dat};
      \addlegendentry{bỏ phiếu trên điểm xáo trộn (trung vị, khoảng 90\%)}
    \end{axis}
  \end{tikzpicture}
  \Description{Hai khung, khung này nằm trên khung kia. Trên: điểm xác thực của 210 tập meta-path của một lần chạy tìm kiếm, sắp từ tốt nhất đến kém nhất, với nhiều tập nằm trong dung sai so với tập tốt nhất. Dưới: các cột cho bốn và sáu meta-path, đếm số tập khác nhau mà 30 lần chạy chọn, cho tập thắng riêng của từng lần chạy, cho bỏ phiếu qua các lần chạy và cho bỏ phiếu trên điểm xáo trộn.}
  \caption{Một lần chạy tìm kiếm không xếp hạng được các tập đứng đầu của nó. (a) Điểm xác thực
    của \val{search-sets} tập mà một lần chạy tính điểm, đã sắp xếp: nhiều tập nằm trong dung sai
    $\delta$ so với tập tốt nhất. (b) Qua \val{search-runs} lần chạy, các tập thắng hầu như đều khác
    nhau. Bỏ phiếu cho mức đồng thuận cao hơn, nhưng chỉ cao hơn đôi chút so với khi bỏ phiếu trên
    điểm xáo trộn.}
  \label{fig:ties}
\end{figure}


\paragraph{Điểm trung bình.} Các lần chạy tính điểm những tập khác nhau, vì nhóm ứng viên cuối của
chúng khác nhau, nhưng xếp hạng giống nhau những tập mà chúng cùng có (hệ số Spearman trung vị
\val{search-agree-4} với $M = 4$ và \val{search-agree-6} với $M = 6$).
Chúng tôi khớp giá trị mất mát xác thực của tập $s$ trong lần chạy $r$ dưới dạng $a_r + b_s$ bằng
bình phương tối thiểu và so sánh \val{search-avg-sets-4} và \val{search-avg-sets-6} tập được tính
điểm trong ít nhất năm lần chạy. Sau đó, chúng tôi áp dụng quy tắc một sai số chuẩn: các tập có $b_s$
cách giá trị của tập tốt nhất không quá một sai số chuẩn (\val{search-se-4} với $M = 4$ meta-path và
\val{search-se-6} với $M = 6$ meta-path) tạo thành một cửa sổ, trong đó tập đơn giản nhất theo tổng
số bước được ưu tiên, vì các meta-path ngắn hơn cũng tạo ra đồ thị HAN thưa hơn. Với $M = 4$
meta-path, cửa sổ này chứa \val{search-window-4} tập: \code{a1} (tập tốt nhất) và \code{a2}. Cả hai
đều được chuyển cho các mô hình. Với $M = 6$ meta-path, nó chứa \val{search-window-6} tập gần ngang
nhau, một vùng phẳng mà tìm kiếm không xếp hạng được. Chúng tôi lấy hai tập đơn giản nhất, bằng
số tập với bốn meta-path. Chúng tôi cũng đã ghi lại lý do cho quyết định này.

\paragraph{Bỏ phiếu.} Một quy tắc thứ hai cho mỗi lần chạy chọn, từ nhóm đồng hạng của nó, tập mà
các meta-path của nó xuất hiện thường xuyên nhất trong các tập thắng của những lần chạy khác. Tập
đồng thuận \code{c} gồm $M$ meta-path phổ biến nhất trong các lựa chọn. Vì bản thân cách bỏ phiếu
đã tạo ra một mức đồng thuận nhất định, chúng tôi so sánh nó với \val{search-null-fleets} loạt chạy
có điểm bị xáo trộn trong từng lần chạy (Hình~\ref{fig:ties}b). Mức đồng
thuận giữa các lần chạy cao hơn mức ngẫu nhiên, nhưng không nhiều: \val{search-null-share-4}\% và
\val{search-null-share-6}\% số loạt chạy xáo trộn có số lựa chọn khác nhau ít như vậy.

\paragraph{Những gì được chuyển cho các mô hình.} Mỗi kích thước chuyển giao \code{a1}, \code{a2}
và \code{c} (Bảng~\ref{tab:sets}). Tập \code{a1} là đối tượng của phép so sánh chính với các tập thủ
công, đã được cố định từ trước. Ngưỡng dung sai $\delta$ được cố định trước các lần chạy cuối, dựa
trên nhiễu giữa các lần chạy trước đó. Các điểm trung bình không phụ thuộc vào nó. Một $\delta$
nhỏ hơn cũng không làm thay đổi \code{a1} lẫn \code{a2}.

\subsection{Các mô hình dùng meta-path}
\label{sec:models}

\paragraph{HAN\@.} HAN \citep{wang2019heterogeneous} biến mỗi meta-path thành một đồ thị trên
các nút đích, chạy cơ chế chú ý mức nút, tức một tầng chú ý đồ thị trên từng đồ thị rồi hợp nhất
các embedding theo từng meta-path bằng cơ chế chú ý mức ngữ nghĩa, tức một trọng số $\beta$ cho
mỗi meta-path dùng chung cho mọi nút. Chúng tôi giữ cả hai mức và thay đổi ba điểm.
\emph{Đồ thị:} đồ thị sản phẩm của mỗi meta-path giữ \val{han-set-topk} láng giềng mạnh nhất của
mỗi sản phẩm theo số đường đi (walk count) sau mỗi bước, vì nếu không cắt tỉa, các meta-path có hai
bước qua người dùng sẽ nối một sản phẩm với phần lớn kho sản phẩm. Đồ thị là nhị phân, đối xứng và
cho mỗi sản phẩm một cạnh tự nối. \emph{Người dùng:} HAN chỉ mã hóa các sản phẩm. Mỗi
người dùng có một embedding tự do và một cặp được tính điểm bằng công thức~\eqref{eq:score} với BPR.
Các embedding người dùng chiếm \val{han-user-params}~triệu trong số
\val{han-params-min}--\val{han-params-max}~triệu tham số của HAN. \emph{Trọng số meta-path:} tầng
chú ý của mọi meta-path bắt đầu từ cùng trọng số, còn $\beta$ được giữ bằng nhau trong
\val{han-set-beta-warmup} epoch đầu. Nếu các tầng chú ý được khởi tạo ngẫu nhiên khác nhau cho từng meta-path, trọng số mà HAN gán cho
một meta-path sẽ phụ thuộc vào vị trí của nó trong danh sách meta-path. Không có giai đoạn khởi động, một meta-path đã chiếm toàn bộ
trọng số chỉ sau vài epoch.

\paragraph{metapath2vec.} metapath2vec \citep{dong2017metapath2vec} học embedding của nút, không
dùng đặc trưng, bằng một mô hình skip-gram trên các bước đi ngẫu nhiên theo một meta-path.
Chúng tôi đi mọi meta-path của một tập cùng nhau, mỗi meta-path như một khuôn mẫu lặp lại
(\code{PUP} cho P--U--P--U\ldots), bắt đầu từ mọi sản phẩm, rồi tính điểm các embedding đã cố định
bằng công thức~\eqref{eq:score}. Người dùng chỉ được tiếp cận qua các meta-path có \code{U}, nên
một tập không có meta-path nào như vậy sẽ để người dùng không được huấn luyện.

\paragraph{GAT và mô hình phổ biến.} Làm đối chứng không dùng meta-path, một mạng chú ý đồ thị
\citep{velickovic2018graph} được huấn luyện với BPR, chạy trên hợp của mọi quan hệ, bỏ qua kiểu nút
ngoại trừ một phép chiếu đầu vào cho mỗi kiểu: đây là mô hình cơ sở đồng nhất của bộ đánh giá chuẩn
HGB \citep{lv2021hgb}. Người dùng và sản phẩm cùng từ một mạng, nên một người dùng được dựng từ
các sản phẩm mà người đó đã nhận xét. Mỗi bước huấn luyện bỏ các cạnh dương của chính nó khỏi đồ
thị mà nó chạy trên, vì một cạnh giữ lại không bao giờ nằm trong đồ thị. Mô hình phổ biến tính
điểm một sản phẩm bằng số người dùng huấn luyện của nó (bằng điểm thì xếp ngẫu nhiên) và đưa ra
cùng một bảng xếp hạng cho mọi người dùng.
